\documentclass[12pt]{article}
\usepackage[margin=0.75in,top=0.8in,bottom=0.65in]{geometry}
\usepackage{lmodern}
\usepackage{microtype}
\usepackage{fancyhdr}
\usepackage{titlesec}
\usepackage{enumitem}
\usepackage{lastpage}
\usepackage[hidelinks]{hyperref}

\setlength{\parindent}{0pt}
\setlength{\parskip}{0.38em}
\setlength{\headheight}{26pt}
\titleformat{\section}{\large\bfseries\scshape}{}{0pt}{}
\titleformat{\subsection}{\normalsize\bfseries}{}{0pt}{}
\pagestyle{fancy}
\fancyhf{}
\fancyhead[C]{\footnotesize Lit Example Reader \quad Assignment ID: U1L1LE \quad Page \thepage\ of \pageref{LastPage}}
\renewcommand{\headrulewidth}{0.5pt}

\newenvironment{playpassage}{\begin{quote}\small\setlength{\parskip}{0.25em}}{\end{quote}}
\newcommand{\speaker}[1]{\textbf{#1.}}

\begin{document}

\begin{center}
{\LARGE\bfseries Week 1 Lit Example Reader}\par\vspace{0.05em}
{\large\scshape Macbeth: Prophecy, Possibility, and Proof}\par\vspace{0.15em}
\rule{0.78\textwidth}{0.5pt}
\end{center}

\textbf{Purpose:} This reader gives two short Macbeth passages for applying Week 1's logic habit: identify the conclusion, stated reason, hidden assumption, and whether the conclusion follows. Read the passages as first-year college students: slowly, precisely, and with attention to what has actually been proved.

\textbf{What to watch for:} Macbeth does not simply hear information. He begins to reason from it. The question is not whether the witches are interesting, frightening, or dramatically powerful. The logical question is: \textit{What follows from what?}

\section*{Before the Passages: The Week 1 Logic Tool}

Whately's opening lesson is modest but demanding. Logic does not tell Macbeth what he wants, what he fears, or what the future will contain. It asks whether a conclusion has been earned by the reasons offered. For each passage, keep four questions separate:

\begin{enumerate}[leftmargin=1.5em,itemsep=0.2em]
\item \textbf{Conclusion:} What claim is the speaker moving toward?
\item \textbf{Stated reason:} What evidence or support is actually given?
\item \textbf{Hidden assumption:} What unstated idea must be true for the reason to prove the conclusion?
\item \textbf{Follow-test:} Does the conclusion really follow, or does the speaker move too quickly?
\end{enumerate}

Do not reduce the scene to a moral slogan. Macbeth's thought is complicated. A careful reader should be able to say both what tempts him and what the logic does or does not prove.

\section*{Passage 1: Act 1, Scene 3 — After the Witches Speak}

\textbf{Context:} The witches have greeted Macbeth with titles: Thane of Glamis, Thane of Cawdor, and future king. Soon afterward, messengers confirm that Macbeth has indeed been named Thane of Cawdor. Macbeth now has one confirmed prediction and one much larger unfulfilled prediction.

\begin{playpassage}
\speaker{Banquo} But 'tis strange:\par
And oftentimes, to win us to our harm,\par
The instruments of darkness tell us truths,\par
Win us with honest trifles, to betray's\par
In deepest consequence.

\speaker{Macbeth} Two truths are told,\par
As happy prologues to the swelling act\par
Of the imperial theme. \ldots\par
This supernatural soliciting\par
Cannot be ill, cannot be good: if ill,\par
Why hath it given me earnest of success,\par
Commencing in a truth? I am Thane of Cawdor:\par
If good, why do I yield to that suggestion\par
Whose horrid image doth unfix my hair\par
And make my seated heart knock at my ribs,\par
Against the use of nature? Present fears\par
Are less than horrible imaginings.\par
My thought, whose murder yet is but fantastical,\par
Shakes so my single state of man\par
That function is smother'd in surmise,\par
And nothing is but what is not.

\speaker{Macbeth} If chance will have me king, why, chance may crown me,\par
Without my stir.
\end{playpassage}

\subsection*{Reader Notes}

Banquo and Macbeth respond differently to the same facts. Banquo gives you a helpful model for the Week 1 method: a statement may be true and still not prove every conclusion a person wants to draw from it. Macbeth has one confirmed prediction and one much larger possibility before him. As you read, separate the fact he actually has from the future conclusion he is tempted to consider.

\textbf{Reading task:} Mark where Macbeth names evidence, where he moves into possibility, and where his imagination begins adding what the witches did not directly command.

\newpage

\section*{Passage 2: Act 1, Scene 7 — Macbeth Tests the Murder}

\textbf{Context:} Duncan is Macbeth's king, kinsman, and guest. Macbeth is considering whether to murder him. In this soliloquy, Macbeth briefly becomes his own critic. He names several reasons not to act, and then he names the motive pushing him forward.

\begin{playpassage}
\speaker{Macbeth} If it were done when 'tis done, then 'twere well\par
It were done quickly: if the assassination\par
Could trammel up the consequence, and catch\par
With his surcease success; that but this blow\par
Might be the be-all and the end-all here,\par
But here, upon this bank and shoal of time,\par
We'ld jump the life to come. But in these cases\par
We still have judgment here; that we but teach\par
Bloody instructions, which, being taught, return\par
To plague the inventor: this even-handed justice\par
Commends the ingredients of our poison'd chalice\par
To our own lips.

He's here in double trust;\par
First, as I am his kinsman and his subject,\par
Strong both against the deed; then, as his host,\par
Who should against his murderer shut the door,\par
Not bear the knife myself.

Besides, this Duncan\par
Hath borne his faculties so meek, hath been\par
So clear in his great office, that his virtues\par
Will plead like angels, trumpet-tongued, against\par
The deep damnation of his taking-off.

I have no spur\par
To prick the sides of my intent, but only\par
Vaulting ambition, which o'erleaps itself\par
And falls on the other.
\end{playpassage}

\subsection*{Reader Notes}

This passage is unusually useful for logic because Macbeth almost argues against himself. Do not let the poetry hide the structure. Find the claims he makes, the kinds of support he offers, and the point at which he stops giving reasons and names what is pushing him. The worksheet will ask you to decide what his reasons actually prove.

\textbf{Reading task:} Track Macbeth's reasons against the murder. At the final turn, mark the word that names his inward push, but decide for yourself how that word functions logically.

\section*{How to Use These Passages on the Worksheet}

The worksheet will not ask you to summarize the plot. It will ask you to apply Week 1 logic to Macbeth's reasoning. When answering, quote briefly if useful, but do not fill the box with long quotation. Your job is to name the logical structure: claim, reason, hidden assumption, and follow-test.

\textbf{Strong college-level answers} should do three things:

\begin{itemize}[leftmargin=1.5em,itemsep=0.2em]
\item separate Shakespeare's dramatic language from the logical movement of the argument;
\item avoid treating prophecy, fear, ambition, or confidence as proof by themselves;
\item explain the hidden assumption rather than merely naming a theme.
\end{itemize}

\end{document}
